% Ad Familiares 11.20 --- headnote
% ad-familiares-11-20, 24 May 43 BC, Eporedia

\textit{Decimus Brutus to Cicero, from Eporedia on 24 May 43 BC ---
Perseus dateline \textit{Scr. Eporediae ix K. Iun. a. 711 (43)}. Eporedia
is the modern Ivrea, at the foot of the Alps; Brutus has paused on the
road, declining to cross over until he knows what is happening at Rome.
The occasion is a warning brought to him by Labeo Segulius, just returned
from Octavian's camp: that the veterans are speaking dangerously of
Cicero, that Octavian and Brutus have been pointedly left off the board
of ten land-commissioners, and that the young Caesar himself has been
heard repeating a saying attributed to Cicero --- that the young man must
be ``praised, decorated, and lifted up'' --- with the dry remark that he
would not let himself be ``lifted out of the way.''}

\textit{The pun on \textit{tollendum} (``raised'' / ``done away with'') is
Latin's, not the translator's: the three gerundives \textit{laudandum,
ornandum, tollendum} read on the surface as a tricolon of advancement,
but \textit{tollere} also means ``to remove,'' and Brutus's gloss
(\textit{se non esse commissurum ut tolli possit}) shows that Octavian, or
his messenger, has heard the second sense. Brutus himself is sceptical
that Cicero said any such thing, and suspects Labeo of either embroidery
or fabrication. The remainder of the letter is a careful piece of
political counsel: meet the veterans halfway on the land-commission and
on the assignments, and watch your back. The ``lifted up'' saying became
one of the famous pieces of testimony for Cicero's relationship with the
young Caesar in the Augustan tradition --- Velleius Paterculus quotes it
in much the form preserved here --- and is the most-cited single line in
the correspondence of these last weeks of Cicero's life. It is reported
here, in the field, six months before Cicero's killing.}
